% =====================================================================
%  CHƯƠNG 1. GIỚI THIỆU                           | Phụ trách: Cả 2
%  Ngân sách: mục tiêu 5 trang, trần 6,0 trang
%  Nguồn: docs/DCDATN_25410104_NguyenMinhNhat_25410034_LeQuangHoaiDuc.pdf,
%         README.md, src/ml/pipeline/onnx_pipeline.py (mô tả pipeline)
% =====================================================================
\chapter{Giới Thiệu}
\label{ch:gioithieu}

\section{Bối Cảnh và Động Lực}

Nhu cầu quản lý bãi giữ xe tại Việt Nam tăng nhanh theo mật độ phương tiện cá nhân.
Phần lớn bãi xe vẫn vận hành thủ công hoặc bằng thẻ từ: nhân viên kiểm tra thẻ hoặc biển số bằng mắt, ghi chép giấy và thu
phí trực tiếp. Cách làm này tốn nhân lực ở cổng, dễ sai sót khi đối chiếu xe vào ra,
không lưu bằng chứng hình ảnh khi tranh chấp và không tạo ra dữ liệu vận hành để thống
kê lưu lượng, doanh thu hay phát hiện bất thường.

Thị giác máy tính và các mô hình phát hiện đối tượng hiện đại cho phép tự động hóa
khâu nhận diện tại cổng, thay quy trình thẻ từ bằng luồng chụp ảnh, phân tích và ghi
nhận tức thì. Mỗi lượt xe trở thành một bản ghi có thời gian, biển số, loại xe và ảnh
toàn cảnh, mở ra khả năng thống kê và phân tích lưu thông.

Xu hướng đưa suy luận xuống thiết bị biên ngày càng rõ trong các hệ thống bãi xe thông
minh~\autocite{pradhan2025iotparking,arukonda2026smartparking,rani2024ips}. Truyền video
từ camera cổng lên máy chủ hoặc cloud có ba nhược điểm: độ trễ tăng khi đường truyền
nghẽn, chi phí máy chủ và băng thông tăng theo số lượng bãi xe, và ảnh phương tiện cùng
biển số liên tục ra ngoài mạng nội bộ. Nhược điểm cuối đáng lưu ý khi Luật Bảo vệ dữ liệu
cá nhân (hiệu lực 01/01/2026) xếp biển số và ảnh nhận dạng vào dữ liệu cá nhân cần kiểm
soát truy cập và hạn chế lưu trữ~\autocite{vn2025lutbvdlcn}. Suy luận tại biên xử lý ngay
tại chỗ, chỉ ghi thông tin thiết yếu vào cơ sở dữ liệu, không truyền video thô ra ngoài.

Các kỹ thuật tối ưu như lượng tử hóa (quantization) và chuyển sang định dạng ONNX đã
thu hẹp khoảng cách hiệu năng giữa thiết bị biên và máy
chủ~\autocite{mdpi2026sbcbenchmark,sonnara2025lightedge}. Raspberry Pi 5 với CPU ARM
Cortex-A76 đủ sức chạy các mô hình phát hiện cỡ nano ở tốc độ chấp nhận được sau khi
xuất ONNX. Đề tài đặt trọng tâm vào việc chạy trọn vẹn một hệ thống quản lý bãi xe
trên một thiết bị biên đơn lẻ, thay vì chỉ dừng ở khâu nhận dạng.

\section{Vấn Đề Đặt Ra}

\textbf{Đặc thù biển số Việt Nam.}
Ô tô dùng biển một hàng, xe máy dùng biển hai hàng, đòi hỏi bước tách dòng trước khi
đọc ký tự; biển hai hàng gần như không xuất hiện trong dữ liệu quốc tế nên các pipeline
ALPR thông dụng xử lý sai hoặc bỏ sót. Seri MĐ và TĐ của xe máy điện lại chứa ký tự
``Đ'' nằm ngoài bảng ký tự của các engine OCR tiếng
Anh~\autocite{tran2025vietnamlpr,le2023motorcycleplate}, nên không thể tái sử dụng trực
tiếp mô hình huấn luyện cho biển số nước ngoài.

\textbf{Thiếu dữ liệu phù hợp.}
Chưa có bộ dữ liệu bãi xe Việt Nam nào gán nhãn đồng thời cho biển số, loại xe và màu
nền biển; các bộ quốc tế như CCPD hay UFPR-ALPR chỉ có vùng biển và ký tự. Phần lớn ảnh
công khai trong nước chụp ở điều kiện ánh sáng tốt, không đại diện cho cổng bãi xe thực
tế với ngược sáng, mưa hay đêm tối. Ghép nhiều nguồn không đồng nhất đòi hỏi làm sạch,
chuẩn hóa nhãn và kiểm soát rò rỉ giữa tập huấn luyện và kiểm tra.

\textbf{Ràng buộc tài nguyên thiết bị biên.}
Thiết bị biên có RAM hạn chế và không có GPU rời, buộc mô hình phải nhỏ và được tối ưu
mới đạt độ trễ chấp nhận được~\autocite{mdpi2026sbcbenchmark,sonnara2025lightedge}.
Pipeline gọi năm bước suy luận liên tiếp trên cùng một khung ảnh nên độ trễ tích lũy
theo chiều dài pipeline; đạt mục tiêu dưới hai giây trên Raspberry Pi 5 đòi hỏi chọn cẩn
thận kích thước mô hình, độ phân giải đầu vào và định dạng tối ưu cho từng bước.

\textbf{Thiếu hệ thống end-to-end có nghiệp vụ.}
Phần lớn công trình dừng ở nhận dạng biển số, không xử lý phiên gửi xe, đối chiếu chống
gian lận và tính phí, nên khoảng cách tới một sản phẩm vận hành được vẫn còn lớn. Một
bãi xe thật cần phân biệt xe vào và xe ra trên cùng biển số, phát hiện biển số vào không
khớp biển số ra (xe hoán đổi biển), áp biểu phí theo nhóm phương tiện, lưu ảnh toàn cảnh
làm bằng chứng và có giao diện tra cứu thời gian thực.

\section{Câu Hỏi Nghiên Cứu}
\label{sec:cauhoi}

Từ các vấn đề trên, câu hỏi nghiên cứu trung tâm được phát biểu như sau: \emph{làm thế
nào để xây dựng một hệ thống quản lý bãi giữ xe tự động hóa khâu nhận diện phương tiện
Việt Nam, vận hành hoàn chỉnh trên thiết bị biên giá rẻ chỉ dùng CPU, đạt độ chính xác
và tốc độ đáp ứng yêu cầu thực tế, đồng thời tuân thủ quy định bảo vệ dữ liệu cá nhân?}
Câu hỏi này được cụ thể hóa thành bốn câu hỏi thành phần:

\begin{enumerate}[leftmargin=*, label=\textbf{CH\arabic*.}]
  \item \textbf{Độ tin cậy.} Hệ thống xác định thông tin phương tiện (biển số, loại xe) từ ảnh cổng chính xác đến đâu, và các thuật toán nghiệp vụ (đối chiếu phiên, tính phí) xử lý thế nào khi mô hình nhận dạng sai?
  \item \textbf{Tối ưu tài nguyên.} Trên phần cứng chỉ có CPU, cần cân bằng độ chính xác, độ trễ và điện năng ra sao để mỗi lượt xe xử lý dưới hai giây?
  \item \textbf{Phương pháp dữ liệu.} Khi chưa thu được dữ liệu thực địa, việc kết hợp các bộ dữ liệu công khai có đủ đại diện để huấn luyện và đánh giá không, và cần kiểm soát gì để tránh đánh giá thiên lệch?
  \item \textbf{Bảo mật và riêng tư.} Cơ chế lưu trữ và tra cứu biển số, ảnh phương tiện cần thiết kế thế nào để không rò rỉ dữ liệu cá nhân, đáp ứng Luật Bảo vệ dữ liệu cá nhân?
\end{enumerate}

\section{Mục Tiêu, Phạm Vi và Tiêu Chí Thành Công}

\subsection{Mục Tiêu Tổng Quát}

Mục tiêu tổng quát là xây dựng một hệ thống quản lý bãi giữ xe thông minh, tự động nhận
diện biển số, loại xe và màu nền biển, thiết kế độc lập nền tảng và ưu tiên triển khai
trên thiết bị biên. Hệ thống gồm cả phần suy luận thị giác máy tính lẫn nghiệp vụ quản
lý vào ra, tính phí và dashboard thống kê, chạy được trên máy chủ, PC và thiết bị biên
với cùng một bộ mã nguồn.

\subsection{Mục Tiêu Cụ Thể}

Năm mục tiêu cụ thể theo đề cương:

\textbf{Mục tiêu 1: Mô đun nhận diện biển số.}
Phát hiện vùng biển số và đọc ký tự (OCR) cho cả biển một hàng của ô tô và biển hai
hàng của xe máy, xử lý đúng ký tự ``Đ'' trong seri biển xe máy điện.

\textbf{Mục tiêu 2: Phân loại phương tiện và nhận màu nền biển.}
Phân loại phương tiện theo bốn lớp (ô tô, xe máy, xe tải, xe buýt) và nhận màu nền biển
(trắng, vàng, xanh, đỏ) để phân nhóm xe cá nhân, kinh doanh và cơ quan nhà nước theo
quy định về biển số của Bộ Công an.

\textbf{Mục tiêu 3: Cơ sở dữ liệu và logic quản lý vào ra.}
Ghi nhận thời gian vào ra, đối chiếu biển số để phát hiện gian lận, lưu ảnh toàn cảnh
làm bằng chứng và tính phí tự động theo nhóm phương tiện.

\textbf{Mục tiêu 4: Dashboard quản lý và thống kê.}
Theo dõi xe ra vào theo thời gian thực, tra cứu phương tiện trong bãi, thống kê theo
loại xe, màu biển, lưu lượng theo giờ và doanh thu.

\textbf{Mục tiêu 5: Thiết kế khả chuyển và tối ưu cho thiết bị biên.}
Thiết kế kiến trúc độc lập nền tảng, tối ưu và đóng gói mô hình bằng ONNX, minh họa khả
năng chạy trên thiết bị biên bằng Raspberry Pi 5.

\subsection{Đối Tượng và Phạm Vi Nghiên Cứu}

Đối tượng nghiên cứu gồm ba nhóm: biển số và phương tiện Việt Nam tại cổng bãi xe (biển
một hàng của ô tô, biển hai hàng của xe máy, bốn màu nền biển); các mô hình thị giác máy
tính cho phát hiện, đọc ký tự và phân loại đủ nhẹ để chạy trên thiết bị biên; và nghiệp
vụ quản lý bãi xe.

Đề tài giới hạn trong mô hình cổng bãi xe, tức phương tiện dừng tại cổng để xử lý thay
vì di chuyển tốc độ cao (free-flow). Phương tiện trọng tâm là xe máy và ô tô, chiếm tỷ
trọng áp đảo tại các bãi xe đô thị Việt Nam; mô hình phân loại vẫn hướng tới đủ bốn lớp
nhưng đánh giá sâu tập trung vào hai lớp chính. Đề tài không xử lý nhận diện dòng xe tốc
độ cao, chưa phân tích nhiều camera đồng thời và không tích hợp kiểm soát cơ học như
barrier hay đèn tín hiệu.

Dữ liệu huấn luyện kết hợp các bộ công khai có nguồn gốc Việt Nam với phần tự gán nhãn
bổ sung. Hệ thống được minh họa trên một thiết bị biên duy nhất là Raspberry Pi 5.

\subsection{Tiêu Chí Thành Công}

Hai tiêu chí thành công định lượng theo cam kết đề cương:

\begin{itemize}
  \item Độ chính xác đọc biển số đạt trên \SI{90}{\percent} trong điều kiện tiêu chuẩn
        (cổng bãi xe, ánh sáng đủ, xe dừng hoặc tốc độ thấp).
  \item Độ trễ end-to-end dưới \SI{2}{giây} mỗi lượt xe khi chạy trên Raspberry Pi 5.
\end{itemize}

\section{Đóng Góp của Đề Tài}
\label{sec:donggop}

\textbf{Pipeline nhận diện năm thành phần;}
phát hiện xe, phân loại loại xe và kiểu dáng, phát hiện biển số, nhận màu nền biển và
đọc ký tự, chạy thành một luồng trên cùng một khung ảnh (Hình~\ref{fig:pipeline}); cùng
một bộ mã suy luận chạy trên cả máy tính lẫn thiết bị biên.

\textbf{Xử lý ký tự ``Đ''.}
Đề tài mở bảng ký tự lên 37 ký tự và sửa các điểm trong tiền xử lý cùng suy luận từng
làm mất ký tự này, nhờ đó đọc đúng biển seri MĐ và TĐ của xe máy điện. Đây là điểm khác
biệt so với các engine OCR tiếng Anh được dùng lại trực tiếp mà không điều chỉnh bảng
ký tự.

\textbf{Hệ thống quản lý bãi xe.}
Hệ thống không chỉ nhận dạng mà còn có phiên gửi xe, đối chiếu chống gian lận, tính phí
theo nhóm xe, phân quyền người dùng và tự động xóa dữ liệu theo Luật Bảo vệ dữ liệu cá
nhân, điều mà phần lớn công trình học thuật về ALPR không đề cập.

% Nguồn hình: pipeline-example.png minh họa trên một ảnh camera cổng thật;
% sơ đồ khối tương ứng ở figures/pipeline.tex (dựng theo onnx_pipeline.py run()).
\begin{figure}[htbp]\centering
\includegraphics[width=0.5\textwidth]{pipeline-example.png}
\caption{Pipeline nhận diện minh họa trên một ảnh camera cổng thật: nhánh biển số (phát hiện biển, nhận màu nền, đọc ký tự) và nhánh xe (định vị vùng xe, phân loại loại xe bằng ResNet18, phân loại kiểu dáng bằng MobileNetV3-Small).}

\label{fig:pipeline}
\end{figure}

\section{Cấu Trúc Báo Cáo}

Chương~\ref{ch:lienquan} khảo sát các công trình liên quan và định vị khoảng trống của
đề tài. Chương~\ref{ch:dulieu} trình bày quy định biển số và dữ liệu.
Chương~\ref{ch:phathien}, \ref{ch:ocr} và \ref{ch:phanloai} lần lượt trình bày phát hiện
biển số, nhận dạng ký tự cùng màu nền biển, và phân loại phương tiện.
Chương~\ref{ch:hethong} mô tả hệ thống quản lý bãi xe, Chương~\ref{ch:trienkhai} đánh
giá triển khai trên Raspberry Pi 5, Chương~\ref{ch:ketluan} tổng kết và nêu hướng phát
triển.
